% !TeX root = ../../../main.tex
\section{The attempted induction}
\label{sec:as_induction_preliminaries}

Let $G$ be a finite abelian group, let $H\subseteq G$ be a subgroup, and let
$\cP\to X$ be a $G$-torsor over an $S$-scheme $X$ that is Zariski-locally
quasi-projective over $S$. Recall the canonical factorization
\[
\cP\longrightarrow\cP/H\longrightarrow X,
\]
where $\cP/H=(G/H)\times^G\cP$ is a $G/H$-torsor over $X$ and
$\cP\to\cP/H$ is an $H$-torsor.

In addition to the kernel $\KG=\ker(G^d\to G)$ of the summation morphism,
put
\[
K_H:=\ker(H^d\to H),\qquad
K_{G/H}:=\ker((G/H)^d\to G/H).
\]
The snake lemma applied to the two summation maps gives an exact sequence
\begin{equation}\label{eq:summation_kernels_ses}
0\longrightarrow K_H\longrightarrow\KG\longrightarrow K_{G/H}
\longrightarrow0.
\end{equation}

\begin{construction}[The symmetric-power comparison morphism]
\label{construction:symmetric_power_comparison}
For the $G$-torsor $\cP\to X$ fixed above, put
\[
\Xi_G:=\KG\rtimes S_d.
\]
The inclusion $S_d\hookrightarrow\Xi_G$, given by
$\sigma\mapsto((1,\ldots,1),\sigma)$, induces a canonical comparison
morphism
\[
q_{\cP}:\cP^{(d)}
=S_d\backslash\cP^{\times_S d}
\longrightarrow
\Xi_G\backslash\cP^{\times_S d}
=\cP^{[d]}.
\]
Indeed, the quotient map from $\cP^{\times_S d}$ to its $\Xi_G$-quotient
is $S_d$-invariant and therefore factors uniquely through the
$S_d$-quotient.
\end{construction}

We use the following four quotient schemes:
\begin{align*}
  \cP^{[d]_G}&:=(\KG\rtimes S_d)\backslash\cP^{\times_S d},\\
  \cP^{[d]_H}&:=(K_H\rtimes S_d)\backslash\cP^{\times_S d},\\
  (\cP/H)^{(d)}&:=(H^d\rtimes S_d)\backslash\cP^{\times_S d},\\
  (\cP/H)^{[d]}&:=((H^d+\KG)\rtimes S_d)
      \backslash\cP^{\times_S d}.
\end{align*}
The first is the symmetric $G$-power of $\cP$. The second is the symmetric
$H$-power relative to $\cP\to\cP/H$, the third is the symmetric power of the
scheme $\cP/H$, and the fourth is the symmetric $G/H$-power of the
$G/H$-torsor $\cP/H\to X$. For the last two identifications, use
$(\cP/H)^{\times_S d}\cong H^d\backslash\cP^{\times_S d}$ and note that the
preimage of $K_{G/H}$ under $G^d\twoheadrightarrow(G/H)^d$ is
$H^d+\KG$.

\begin{prop}\label{prop:torsor_tower_diagram}
The four quotient schemes fit into a canonical commutative diagram
\[
\begin{tikzcd}[row sep=large, column sep=large]
  \cP^{[d]_H}\arrow[r,"H"]\arrow[d]
    &(\cP/H)^{(d)}\arrow[d,"{q_{\cP/H}}"]\\
  \cP^{[d]_G}\arrow[r,"H"']\arrow[dr,"G"',bend right=20]
    &(\cP/H)^{[d]}\arrow[d,"G/H"]\\
    &X^{(d)},
\end{tikzcd}
\]
in which every labelled arrow is a torsor under the indicated constant
group. The upper square is Cartesian:
\[
\cP^{[d]_H}\cong
\cP^{[d]_G}\times_{(\cP/H)^{[d]}}(\cP/H)^{(d)}.
\]
The vertical map $q_{\cP/H}$ is the comparison morphism of
\Cref{construction:symmetric_power_comparison}, applied to the
$G/H$-torsor $\cP/H\to X$.
\end{prop}

\begin{proof}
Apply \Cref{prop:symmetric_power_torsor} to the $G$-torsor $\cP\to X$, the
$H$-torsor $\cP\to\cP/H$, and the $G/H$-torsor $\cP/H\to X$. This gives the
three labelled torsors; the comparison map $q_{\cP/H}$ is supplied by
\Cref{construction:symmetric_power_comparison}, applied to
$\cP/H\to X$. The factorization
\[
\cP^{[d]_G}\longrightarrow(\cP/H)^{[d]}\longrightarrow X^{(d)}
\]
is the canonical quotient decomposition of the $G$-torsor
$\cP^{[d]_G}\to X^{(d)}$ by $H$, because both its intermediate quotient and
$(\cP/H)^{[d]}$ are realized as
\[
((H^d+\KG)\rtimes S_d)\backslash\cP^{\times_S d}.
\]

Let
\[
T:=\cP^{[d]_G}\times_{(\cP/H)^{[d]}}(\cP/H)^{(d)}.
\]
Base change makes $T$ an $H$-torsor over $(\cP/H)^{(d)}$. The quotient maps
give an $H$-equivariant morphism $\cP^{[d]_H}\to T$. A morphism of
$H$-torsors is an isomorphism, proving the Cartesian assertion.
\end{proof}

\begin{rem*}
The two vertical comparison maps in the diagram need not be torsors under a
constant group. For example, $K_H\rtimes S_d$ need not be normal in
$\KG\rtimes S_d$, because coordinate permutations act nontrivially on
$\KG/K_H$. Thus the map $q_{\cP/H}$ is not another torsor layer to which the cyclic
induction hypotheses automatically apply. Its ramification at the boundary
is precisely the extra datum that will matter below.
\end{rem*}

We now place this quotient tower in the geometric setting relevant to the
attempted induction.
Let $C$ be a smooth projective curve, let $U=C\setminus\mdls$, and let
$\cP\to U$ be a $G$-torsor. Fix $p_\infty\in\mdls$ and an integer $d\geq1$,
and assume that the ramification of $\cP$ at $p_\infty$ is bounded by $d$ in
the sense of \Cref{definition:local_ramification_boundness}.

The $G/H$-torsor $\cP/H\to U$ compactifies to a finite morphism of smooth
projective curves
\[
f:C'\longrightarrow C.
\]
Choose $p'_\infty\in f^{-1}(p_\infty)$. The induced morphism on symmetric
powers sends $d\cdot p'_\infty$ to $d\cdot p_\infty$:
\[
f^{(d)}:C'^{(d)}\longrightarrow C^{(d)}.
\]
Let
\[
\pi_C:X_C\longrightarrow C^{(d)},\qquad
\pi_{C'}:X_{C'}\longrightarrow C'^{(d)}
\]
be the blowups at these two diagonal points. Write $E_C,E_{C'}$ for their
exceptional divisors, $\eta_C,\eta_{C'}$ for the respective generic points,
and
\[
V_C:=\Oo_{X_C,\eta_C}\subset K(C^{(d)}),\qquad
V_{C'}:=\Oo_{X_{C'},\eta_{C'}}\subset K(C'^{(d)})
\]
for the corresponding DVRs.

The quotient tower, compactifications, and blowups fit into the diagram
\[
\begin{tikzcd}[row sep=2.4em, column sep=2.2em]
\cP^{[d]_H}\ar[r,"H"]\ar[d]
  &(\cP/H)^{(d)}\ar[d,"{q_{\cP/H}}"]\ar[r,hookrightarrow,"\iota'"]
  &C'^{(d)}\ar[dd,"f^{(d)}"]
  &X_{C'}\ar[l,"\pi_{C'}"']\ar[dd,dashed,"f_X"]\\
\cP^{[d]_G}\ar[r,"H"']\ar[dr,"G"',bend right=10]
  &(\cP/H)^{[d]}\ar[d,"G/H"]
  &&\\
  &U^{(d)}\ar[r,hookrightarrow,"\iota"']
  &C^{(d)}
  &X_C\ar[l,"\pi_C"]
\end{tikzcd}
\]
The left block is the torsor tower of
\Cref{prop:torsor_tower_diagram}. The middle block records the open
immersions and compactifications. The right block records the blowups.
{\color{blue}
\noindent\textbf{First obstruction: incomparable exceptional valuations.}
\par}
{\color{red}
The
rational lift $f_X$ is defined at the generic point of the exceptional
divisor and sends $\eta_{C'}$ to $\eta_C$, producing the extension of DVRs
\[
V_C\hookrightarrow V_{C'}.
\]
\par}
{\color{blue}
The rational lift $f_X$ is defined at $\eta_{C'}$ by the properness of
$X_C\to C^{(d)}$, but its image need not be $\eta_C$. Consequently, the
inclusion of function fields need not induce an extension of the two DVRs
$V_C\hookrightarrow V_{C'}$.

For example, let $k$ be algebraically closed of characteristic $2$, let
$C'=\mathbb P^1_t$, $C=\mathbb P^1_x$, and take
\[
x=\frac{t^2}{1-t},\qquad d=3,\qquad p_\infty=p'_\infty=0.
\]
This is a separable Artin--Schreier map, since $y=t^{-1}$ satisfies
$y^2-y=x^{-1}$. Write $a,b,c$ for the elementary symmetric functions
of $t_1,t_2,t_3$, and $s_1,s_2,s_3$ for those of the corresponding
$x$-coordinates. With $D=1+a+b+c$, direct expansion gives
\[
s_1=\frac{a^2+ab+ac+c}{D},\qquad
s_2=\frac{b^2+bc}{D},\qquad
s_3=\frac{c^2}{D}.
\]
At $\eta_{C'}$, the ordinary exceptional valuation has
$v(a)=v(b)=v(c)=1$, so
\[
\bigl(v(s_1),v(s_2),v(s_3)\bigr)=(1,2,2).
\]
Hence the lifted map sends $\eta_{C'}$ to the special point
$[1:0:0]$ of the target exceptional $\mathbb P^2$, rather than its generic
point. In particular, $s_2/s_1$ is a unit at $\eta_C$ but has positive
valuation at $\eta_{C'}$, ruling out the asserted extension of these DVRs.
\par}

We can now formulate the proposed induction step.
In the cyclic $p$-power setting described in the introduction to this
appendix, the two intended inputs are:
\begin{description}
  \item[\textnormal{(IH$_{r-1}$)}]
    $(\cP/H)^{[d]}\to U^{(d)}$ is unramified at $\eta_C$;
  \item[\textnormal{(AS)}]
    $\cP^{[d]_H}\to(\cP/H)^{(d)}$ is unramified at $\eta_{C'}$.
\end{description}
The first is the desired contribution of induction applied to the quotient
$G/H$. The second is the desired contribution of the degree-$p$
Artin--Schreier calculation applied to the $H$-torsor over the compactified
curve $C'$. The proposed induction step asks whether these imply that
\[
\cP^{[d]_G}\longrightarrow U^{(d)}
\]
is unramified at $\eta_C$.
{\color{red}
We will see that the two assertions do not, by
themselves, control the ramification needed for this conclusion.
\par}
{\color{blue}
The valuation mismatch above already prevents a direct comparison. The
next section explains the further obstruction that remains even if one
adds the compatibility hypothesis
$V_{C'}\cap K(C^{(d)})=V_C$.
\par}

